\section{Тестирование}

\subsection{Сборка проекта CMake}

Проект содержит статическую библиотеку \texttt{exam\_grade\_set}, демонстрационную программу \texttt{main} и исполняемый файл \texttt{test\_exam\_grade\_set}. Google Test 1.14.0 подключается через \texttt{FetchContent}, тесты регистрируются в CTest.

\lstinputlisting[style=cmakestyle]{../../CMakeLists.txt}

\subsection{Результаты тестирования}

Реализовано 19 тестов Google Test. Проверяются пустое множество, добавление, повторное добавление, обновление оценки, поиск и удаление по id, сохранение порядка, цепочки операторов, независимость копий, перемещение, самоприсваивание, сравнение, операции множеств и формат вывода. Стресс-тест выполняет 1000 добавлений, обновления оценок, повторные добавления и удаление всех записей.
Все 19 тестов успешно пройдены.
Команды сборки и запуска под Linux из каталога \texttt{lab\_3}:
\begin{verbatim}
cmake -S . -B build
cmake --build build
./build/main
ctest --test-dir build --output-on-failure
\end{verbatim}

\subsection{Демонстрационный пример}

Для первого множества оценка студента 1 обновляется с 3 до 5. Далее выводятся оба множества, их объединение, пересечение и разность:
\begin{verbatim}
{ (1, 5), (2, 4) }
{ (2, 5), (3, 4) }
{ (1, 5), (2, 5), (3, 4) }
{ (2, 4) }
{ (1, 5) }
\end{verbatim}
